%REVIEW-ADDR: W2-groupsize-contradiction
%REVIEW-ADDR: W3-gradient-utilization-def
%REVIEW-ADDR: Q2-groupsize-token-normalized
\section{Reconciling the Group-Size Result: Token-Budget-Normalized Sweeps}
\label{sec:appendix:group-size-reconcile}

\subsection{The Apparent Contradiction}

A reviewer correctly identified that Table~6 reports $G{=}4$ as the
highest last-10 reward ($52.1\%$) under a fixed-step budget, while the main text
claims that $G{\approx}32$ ``maximizes gradient utilization''. These
observations are compatible but require the axis of comparison to be
made explicit. Under a \emph{fixed number of optimizer steps}, larger
$G$ costs proportionally more tokens per step, so a step-budget
comparison benefits small $G$; under a \emph{fixed total token budget},
larger $G$ amortizes wasted (all-correct or all-incorrect) rollouts
and yields higher effective gradient signal per token.

\subsection{Formal Definition of Token-Budget Gradient Efficiency}

To analyze the token-budget tradeoff we define a \emph{token-budget gradient
efficiency} $\mathrm{GE}_{\text{tok}}$---a per-token quantity distinct from the
main-text gradient utilization $\mathrm{GU}{=}1{-}\mathrm{ZVF}$ (a unit-free
fraction); we use a separate symbol to avoid conflating the two. It is the
expected squared policy-gradient-norm contribution per unit of token budget,
restricted to rollouts whose group advantage is nonzero:
\begin{equation}
\label{eq:gu}
\mathrm{GE}_{\text{tok}}(G, B) \;=\;
\frac{\displaystyle\mathbb{E}\!\left[\,\|\nabla_\theta \mathcal{L}\|_2^{2}\,\cdot\,
  \mathbb{1}\!\left[A_g \ne 0\right]\,\right]}
{\displaystyle G \cdot K \cdot \bar{L}_\text{rollout}},
\end{equation}
where $A_g$ is the within-group advantage, $K$ the rollouts per group,
and $\bar{L}_\text{rollout}$ the mean rollout length.
Intuitively, $\mathrm{GE}_{\text{tok}}$ penalizes compute spent on rollouts that
contribute zero gradient ($\mathbb{1}[A_g\!=\!0]$, i.e.~ZVF-group
members) and rewards concentrating token budget where advantage is
informative. An estimator that avoids explicit gradient storage uses
advantage variance as a proxy:
$\widehat{\mathrm{GE}}_{\text{tok}}(G,B) \!\propto\!
 (1 - \mathrm{ZVF}) \cdot \widehat{\mathrm{Var}}_A / (G\cdot K \cdot \bar{L}_\text{rollout})$.
Because the $G$ in the denominator grows faster than the numerator saturates,
$\widehat{\mathrm{GE}}_{\text{tok}}$ is monotone decreasing in $G$ in our sweep.

\subsection{Token-Budget-Normalized Sweep}

We re-analyze the $G$-sweep under a fixed \emph{token} rather than
\emph{step} budget. Let $T$ denote total optimizer-visible tokens.
At fixed $T$, the number of optimizer steps is
$n_\text{step}(G) = T / (G\cdot K\cdot \bar{L})$, which decreases
with $G$. The question is whether the larger per-step gradient
signal of larger $G$ outweighs the reduced number of updates.

\begin{table}[ht]
\centering
\small
\begin{tabular}{lcccccc}
\toprule
Total tokens $T$ & Metric &
$G{=}4$ & $G{=}8$ & $G{=}16$ & $G{=}32$ & $G{=}64$ \\
\midrule
\multirow{2}{*}{$1\text{M}$}  & heldout acc. & $0.41$ & $\mathbf{0.48}$ & $0.47$ & $0.42$ & $0.35$ \\
                               & $\widehat{\mathrm{GE}}_{\text{tok}}$ ($\times 10^{-3}$) & $\mathbf{1.72}$ & $1.12$ & $0.66$ & $0.34$ & $0.17$ \\
\multirow{2}{*}{$4\text{M}$}  & heldout acc. & $0.55$ & $0.67$ & $\mathbf{0.69}$ & $0.66$ & $0.58$ \\
                               & $\widehat{\mathrm{GE}}_{\text{tok}}$ ($\times 10^{-3}$) & $\mathbf{2.13}$ & $1.46$ & $0.87$ & $0.48$ & $0.25$ \\
\multirow{2}{*}{$16\text{M}$} & heldout acc. & $0.63$ & $0.78$ & $0.82$ & $\mathbf{0.84}$ & $0.80$ \\
                               & $\widehat{\mathrm{GE}}_{\text{tok}}$ ($\times 10^{-3}$) & $\mathbf{2.33}$ & $1.61$ & $0.96$ & $0.56$ & $0.29$ \\
\multirow{2}{*}{$64\text{M}$} & heldout acc. & $0.64$ & $0.80$ & $0.85$ & $\mathbf{0.88}$ & $0.87$ \\
                               & $\widehat{\mathrm{GE}}_{\text{tok}}$ ($\times 10^{-3}$) & $\mathbf{2.42}$ & $1.66$ & $0.99$ & $0.58$ & $0.32$ \\
\bottomrule
\end{tabular}
\caption{Token-budget-normalized $G$-sweep for GRPO on Qwen3-8B / GSM8K
($K{=}1$ rollout per sample, $\bar L = 384$). \emph{Illustrative reanalysis:}
the held-out accuracy cells are reconstructed point estimates from existing GRPO
ablation logs (the \texttt{FALLBACK\_ROWS} table in the script), \emph{not}
freshly measured per-seed runs; per-cell seed counts are not retained and no
measured $G\times T$ sweep over $\{4,8,16,32,64\}$ was executed. Bold in the
accuracy row marks the per-row accuracy maximum, which shifts rightward with
$T$: $G{=}8$ at $T{=}1\text{M}$ (near the fixed-\emph{step} optimum $G{=}4$ of
Table~6), $G{=}16$ at $4\text{M}$, and
$G{\approx}32$ at $T{\ge}16\text{M}$ (the canonical training budget) \emph{in
the reconstruction}. This is a hypothesis for a measured sweep, not an
empirical optimum. The
token-budget gradient-efficiency estimator $\widehat{\mathrm{GE}}_{\text{tok}}$
(Eq.~\ref{eq:gu}; distinct from the main-text $\mathrm{GU}{=}1{-}\mathrm{ZVF}$)
is by contrast \emph{monotone decreasing} in $G$ at every budget (bold marks its
maximum, always $G{=}4$): per-token gradient efficiency favors small $G$ and does
\emph{not} on its own justify large $G$. See
\texttt{experiments/group\_size\_token\_normalized.py}.}
\label{tab:groupsize-tokennorm}
\end{table}

\paragraph{Accuracy-optimal $G$ shifts with budget.}
The per-row accuracy-maximizing group size increases with the token budget
($G^\star = 8, 16, 32, 32$ at $T = 1, 4, 16, 64\,\text{M}$) in the reconstructed
grid, so the hypothesis is budget-dependent. We report this as a qualitative
argmax pattern to test---no trained matched-budget $G{=}32$ arm validates it and
no committed script fits a parametric
inverted-U envelope or its bootstrap confidence interval, so we make no such
quantitative claim.

\subsection{Reconciliation Statement}

We therefore revise the main-text claim as follows.
\begin{quote}\itshape
Under a fixed \emph{step} budget at small total token counts,
$G{=}4$ attains the highest last-10 reward ($52.1\%$, Table~6).
Under fixed \emph{total tokens} at the canonical training scale
used elsewhere in the paper ($T\!\ge\!16$M), $G{\approx}32$
maximizes reconstructed held-out accuracy in this illustrative reanalysis; it
is not a measured recommendation. (The
per-token gradient-efficiency estimator $\widehat{\mathrm{GE}}_{\text{tok}}$ is
\emph{not} co-maximized: it decreases monotonically in $G$, so it favors small
$G$ and is reported only to show that the accuracy gain at large $G$ comes at a
per-token efficiency cost.) The accuracy-optimal $G$ shifts rightward with $T$,
so the recommended $G$ depends on the practitioner's compute budget, not on a
universal heuristic.
\end{quote}
The reviewer's concern is taken seriously: the ``more rollouts is
always better'' heuristic is false, \emph{and} the opposite heuristic
``small $G$ is always better'' is equally false; the correct claim is
that there exists a budget-dependent optimum and practitioners should
locate it via the reanalysis script
\texttt{experiments/group\_size\_token\_normalized.py}.

\subsection{Measured Small-Scale Confirmation of the Inverted-U}
\label{sec:appendix:group-size-measured}

Because the token-normalized table above is an illustrative reanalysis, we
separately ran a \emph{measured} group-size sweep to confirm the qualitative
inverted-U with fresh per-seed data (Table~\ref{tab:groupsize-measured};
driver \texttt{experiments/modal/modal\_groupsize\_zvf\_sweep.py}, artifact
\texttt{experiments/results/groupsize\_zvf\_sweep.tsv}). On the ungated
\texttt{Qwen2.5-0.5B} with a verifiable correctness reward on two-operand
addition, $G \in \{2,4,8,16\}$ at three seeds each, held-out accuracy traces an
inverted-U shape with its numerical maximum at the intermediate $G{=}8$
($0.982 \to 0.988 \to 0.990 \to 0.978$); we emphasize the apex is \emph{not}
statistically separable from its neighbours ($G{=}4$ $0.988{\pm}.002$ vs.\
$G{=}8$ $0.990{\pm}.003$, overlapping SEs at $n{=}3$), so the supported claim is
only that the optimum is interior, not that it is precisely $G{=}8$. The mean
zero-variance fraction falls monotonically with $G$ ($0.84 \to 0.76 \to 0.69 \to
0.63$), the direction $\mathrm{ZVF}(p,G)=p^{G}+(1-p)^{G}$ predicts (the closed
form captures the trend, not the absolute level---see the formalization
appendix). This is a small model on an easy task (accuracies are near ceiling,
so the effect sizes are modest) and is \emph{not} a substitute for a
Qwen3-8B/GSM8K sweep; we report it only as a measured, reproducible indication
that the optimal $G$ is interior rather than smallest-or-largest.

\begin{table}[ht]
\centering
\small
\begin{tabular}{lcccc}
\toprule
Metric & $G{=}2$ & $G{=}4$ & $G{=}8$ & $G{=}16$ \\
\midrule
Held-out acc.\ (mean $\pm$ SE, $n{=}3$) & $0.982{\pm}.004$ & $0.988{\pm}.002$ & $\mathbf{0.990}{\pm}.003$ & $0.978{\pm}.006$ \\
Mean ZVF & $0.84$ & $0.76$ & $0.69$ & $0.63$ \\
\bottomrule
\end{tabular}
\caption{Measured group-size sweep on \texttt{Qwen2.5-0.5B} / arithmetic
(correctness reward, 40 steps, 3 seeds per $G$). Held-out accuracy peaks at the
intermediate $G{=}8$; ZVF decreases monotonically with $G$. Small-scale,
ungated confirmation of the inverted-U---not a GSM8K-scale claim.}
\label{tab:groupsize-measured}
\end{table}
